\chapter{Grunnen}

\begin{motto}
Form er ein posisjon i eit rom av moglegheiter,\\ forma av samtidige seleksjonstrykk, navigert av agentar på fleire skalaer.
\end{motto}

Eg lova inga kur som mildnar diagnosen, og eg held ordet. Dette kapitlet gjev ikkje faget tilbake det dei føregåande kapitla tok frå det. Det gjer noko strengare: det viser at det finst ein grunn å byggje på, og difor at faget ikkje har orsakinga si lenger. Så lenge ein trur at design er konstitusjonelt utan fundament — at problema er for vrange, kunnskapen for taus, forma for fin til å fangast — kan ein bere fråvêret av eit fundament som ein tragisk naudsynt tilstand. I det augneblinken det vert vist at eit fundament er mogleg, vert fråvêret av eitt eit val. Diagnosen i kapittel ein til åtte var: faget har ikkje eit fundament. Tillegget i dette kapitlet er verre for faget: det \emph{kunne} hatt eitt heile tida.

Eg skal ikkje leggje fram fundamentet i sin fulle form her; det er det traktaten \textsc{Formlære} gjer, og denne boka er hennar destruktive tvilling, ikkje hennar samandrag. Men eg skuldar lesaren å vise nok til at påstanden ikkje står naken: at det finst ein presis, falsifiserbar modell av nett det faget har sagt ikkje let seg modellere.

\section{Éi setning}

Heile rammeverket kan samlast i éi setning. Form er ein posisjon i eit rom av moglegheiter, forma av samtidige seleksjonstrykk, navigert av agentar på fleire skalaer. Les henne sakte, for kvart ledd gjer arbeid faget aldri har fått gjort.

\emph{Eit rom av moglegheiter.} Forma er ikkje ein ting med eigenskapar; ho er eit punkt i eit formrom der kvar akse er ein målbar dimensjon — høgd, krumming, masse, ornamenttettleik, kva som helst som varierer. Med ein gong forma er eit punkt i eit rom, kan vi snakke presist om avstand, retning, naboskap, variasjon. Stilhistoria sin katalog vert til ei geometri.

\emph{Samtidige seleksjonstrykk.} Det som verkar på forma er ikkje «funksjonen» i eintal, men eit knippe gradientar — ergonomi, materialøkonomi, kulturell meining, nedbryting, makt — som kvar dreg forma i si retning. Forma som vert til er likevekta mellom dei. Her ligg svaret på den tomme formelen frå kapittel \textsc{ii}: forma følgjer ikkje éin funksjon, ho er resten etter mange samtidige trykk, og difor er underbestemminga ikkje eit mysterium men sjølve strukturen. Og her ligg svaret på det vrange problemet frå kapittel \textsc{v}: at trykka er mange og motstridande er ikkje ein grunn til at teori er umogleg; det er innhaldet i teorien. Biologien modellerte nett denne situasjonen for nitti år sidan, med tilpassingslandskapet\kref{12}.

\emph{Navigert av agentar.} Den som gjev form er ein agent som navigerer landskapet med ein avgrensa horisont — ho kan representere nokre aksar og er blind for andre. Her ligg svaret på kapittel \textsc{vi}: svikten til ei form er leseleg som signaturen av kva aksar agenten ikkje representerte. Det manglande nei-et får eit presist uttrykk: det du ikkje kan spesifisere, kan du ikkje forsvare.

\emph{På fleire skalaer.} Agenten er ikkje åleine. Celle, materiale, verktøy, handverkar, organisasjon, marknad — alle er agentar med eigne horisontar som formar landskapet for kvarandre samstundes. Designaren er ikkje den einerådande; han er éin agent blant mange, og det er nett det som gjer at etikken fell ut av modellen i staden for å boltast på, slik kapittel \textsc{vi} kravde.

\section{Éi likning}

Setninga har ei formell form. Den generative modellen i traktaten skriv forma som

\begin{center}
\(\mathrm{Form} = SG\!\left(\msym{∇}_{C(A)}\, L(c,t)\right) \,\msym{∩}\, \pic{◇}\)
\end{center}

der \(SG\) er formgrammatikken\kref{9} som genererer kva former som er moglege, \(L(c,t)\) er tilpassingslandskapet\kref{12} av samtidige trykk, \(\msym{∇}_{C(A)}\) er gradienten slik agenten \(A\) med sin kunnskap \(C\) kan estimere han, og \(\pic{◇}\) er den kognitive horisonten\kref{13} som avgrensar kva agenten kan representere. Tre uavhengige tradisjonar — formgrammatikken til Stiny\kref{9}, tanke--kunnskapsteorien til Hatchuel og Weil\kref{10}, og det substrat-uavhengige agensrammeverket til Levin og Fields\kref{13} — møtest i dette uttrykket. At tre tradisjonar som aldri snakka saman konvergerer på same struktur, er sjølv eit teikn på at strukturen grip noko reelt.

Eg legg ikkje fram likninga for å overtyde deg om at ho er rett; det krev traktaten, fundamentet og falsifiseringsvilkåra. Eg legg henne fram for å vise éi ting, og berre éi: at det \emph{lèt seg gjere}. At det er mogleg å skrive ned, presist og falsifiserbart, kva ei form er og kva som verkar på henne. I det augneblinken er faget si djupaste sjølvforståing — at slik presisjon er prinsipielt umogleg for design — motbevist ved eit eksempel. Eitt fungerande fundament er nok til å gjere fråvêret av eitt til eit val.

\section{Kva ei utdanning på grunn ville krevje}

Eg skal ikkje teikne ein ny studieplan; det ville vere å lyge om kor lett overgangen er. Men retninga følgjer av diagnosen, ledd for ledd, og er verdt å seie reint.

Ei utdanning på grunn ville krevje skriftleg grunngjeving for kvart formval: kva aksar vart representerte, kva vart valde bort, kven ber kostnaden. Ikkje fordi byråkrati er bra, men fordi eit val som ikkje kan grunngjevast slik, ikkje er eit designval; det er eit innfall. Ho ville krevje tid i substratet, ikkje berre i representasjonen — at studenten har arbeidd i materialet, stått i svikten, sett kva forma gjer over tid. Ho ville krevje konsultasjon med nabofaga der forma rører levande system, store menneskemengder, økologi, makt — ikkje som eit påbolta etikkemne, men fordi modellen gjer designaren til éin agent blant mange. Og ho ville krevje eit organ som kan utelukke for brot, slik kvart ekte profesjonsfag har; for utan ei grense å bli utelukka frå, finst det inga profesjon, berre eit yrke. Ein kandidat som aldri har sagt nei til eit oppdrag, har ikkje fullført utdanninga; ho har fullført ei kopiering.

Sjå at kvart av desse krava er den direkte motsetnaden til ein veikskap eg har skildra. Mot blikket som ikkje kan grunngje seg sjølv: kravet om skriftleg grunngjeving. Mot pensumet av udigererte lån: kravet om tid i substratet. Mot etikken som vart bolta på: konsultasjonen som fell ut av modellen. Mot faget som ikkje kan seie nei: organet som kan utelukke. Diagnosen og kuren er same liste, lesen to vegar.

\vignett

Det er alt eg vil seie om grunnen her, for resten står andre stader. Det som står att er ikkje eit argument, men eit val — faget sitt val, og kvar einskild sitt. Det er etterordet.
